%!TEX root = ./main.tex

\section[Routing]{Model Routing}

% SECTION TIMING: 45:00--48:00 (3 frames): 70 + 70 + 40 s.

% Teaching note (70 s): the idea is almost social: a receptionist who answers the
% easy questions personally and forwards the hard ones to a specialist. Per query,
% a router sends easy requests to the local small model and hard ones to the cloud.
% Most requests are easy (set a timer, summarize a note), so the expensive path is
% exercised rarely -- the average cost drops while the worst-case quality is kept.
\begin{frame}{Trick 5: send the question, not every question}
  \vspace{0.7em}
  \centering
  \begin{tikzpicture}[scale=0.85,transform shape]
    \node[endnode,minimum width=1.9cm] (q) at (0,0) {request};
    \node[nnnode,minimum width=1.9cm,minimum height=0.9cm] (r) at (3.0,0)
      {Router};
    \node[nnnode,minimum width=3.3cm,minimum height=1.0cm] (loc) at (7.4,1.1)
      {Local small model\\[-2pt]{\scriptsize\mdseries ``set a timer'',
        summaries}};
    \node[chnode,minimum width=3.3cm,minimum height=1.0cm] (cl) at (7.4,-1.1)
      {Cloud big model\\[-2pt]{\scriptsize\mdseries the genuinely hard asks}};
    \draw[flow] (q) -- (r);
    \draw[flow,draw=softgreen,thick] (r) -- node[lbl,above left] {most} (loc);
    \draw[flow,draw=coreorange] (r) -- node[lbl,below left] {few} (cl);
  \end{tikzpicture}

  \vspace{0.8em}
  \takeawayat{2-}{Average cost of the small model,\\
    worst-case quality of the big one.}
\end{frame}

% Teaching note (70 s): how does the router know a request is hard, before
% answering it? Three practical signals, cheapest first: (1) task type and length
% -- a timer is local by rule, a 30-page analysis is cloud by rule; (2) the small
% model's own confidence -- answer locally, and escalate when its token
% probabilities look unsure (students saw confidence in the PyTorch lecture);
% (3) a small learned router trained on pairs where people preferred the big
% model's answer. This is the shipped architecture: Apple's requests route between
% the ~3B on-device model and Private Cloud Compute; Pixel routes between Gemini
% Nano and the cloud Gemini.
\begin{frame}{How the router decides}
  \vspace{0.6em}
  \raggedright
  \begin{enumerate}[<+->]\itemsep6pt
    \item \textbf{Rules}: task type and length --- timers local,
      30-page analysis cloud
    \item \textbf{Confidence}: answer locally; escalate when the small model's
      own token probabilities look unsure
    \item \textbf{A learned router}: trained on requests where people preferred
      the big model's answer
  \end{enumerate}

  \vspace{0.7em}
  \uncover<4->{\raggedright This is the shipped design: Apple routes between its
    on-device model and Private Cloud Compute; Pixel between Gemini Nano and
    cloud Gemini.\par}
  \plainsource{Apple, ``Introducing Apple's On-Device and Server Foundation
    Models,'' 2024.}
\end{frame}

% Teaching note (40 s): close the five-trick arc with the complete picture: the
% modern assistant is not one model, it is a routed system of models -- quantized
% local ones, MoE cloud ones, caches and split points in between. Hand to wrap-up.
\begin{frame}{The assistant was never one model}
  \vspace{1.0em}
  \raggedright
  What answers you is a \textbf{system}: a router, a quantized on-device model,
  a KV cache on a diet, and an MoE in the cloud --- with split points and
  drafts in between.\par

  \vspace{1.0em}
  \takeawayat{2-}{``Where should the model run?'' ---\\
    the shipped answer is: \textbf{everywhere, each part where it is cheapest.}}
\end{frame}
